% Same geometry as the strategy slide, focusing on quantum preparation.
% Arguments: vertical position, verifier label, arrow label.
\newcommand{\preparesendprotocol}[3]{%
  \begin{scope}[shift={(0,#1)},scale=0.68,transform shape]
    \node[party] at (-4.15,0) {#2};
    \node[party] at (4.15,0) {Prover};
    \draw[communication] (-2.35,0) --
      node[above=7pt,align=center,font=\large,text=CVQCInk] {#3}
      (2.35,0);
  \end{scope}%
}

% outline.md, lines 47--52.
\begin{frame}{CVQC Strategy: Prepare-and-Send}
  \centering
  \begin{tikzpicture}[
      >={Latex[length=3mm]},
      party/.style={circle,draw=CVQCInk,fill=CVQCFog,very thick,
        minimum size=2.5cm,text width=2.1cm,inner sep=0pt,
        align=center,font=\large\bfseries,text=CVQCInk},
      communication/.style={->,line width=1.5pt,draw=CVQCTeal},
      protocolbox/.style={draw=CVQCInk,thick,rounded corners=4mm},
      modellabel/.style={fill=white,inner xsep=6pt,inner ysep=2pt,
        font=\bfseries,text=CVQCInk}
    ]
    \path[use as bounding box] (-6.5,-3.15) rectangle (6.5,3.15);

    \draw[protocolbox] (-5.4,0.4) rectangle (5.4,2.9);
    \node[modellabel] at (-4.25,2.9) {$\mathsf{QPIP}_1$};
    \preparesendprotocol{1.65}{Weak\\quantum\\verifier}%
      {$\ket{+_{\theta}}=T^{\theta}\ket{+}$}

    \only<2>{
      \draw[->,line width=1.5pt,draw=CVQCTeal] (0,0.30) -- (0,-0.30)
        node[midway,right=8pt,font=\small\bfseries,text=CVQCTeal] {compile};

      \draw[protocolbox] (-5.4,-2.9) rectangle (5.4,-0.4);
      \node[modellabel] at (-4.25,-0.4) {$\mathsf{QPIP}_0$};
      \preparesendprotocol{-1.65}{Classical\\verifier}{Remote state\\preparation}
    }
  \end{tikzpicture}

  \note<1>{
    \footnotesize
    In the prepare-and-send paradigm, the weak quantum verifier prepares
    single-qubit states and sends them to the prover. For example,
    $\ket{+_{\theta}}=T^{\theta}\ket{+}
    =(\ket{0}+e^{i\pi\theta/4}\ket{1})/\sqrt{2}$, where
    $T=\operatorname{diag}(1,e^{i\pi/4})$ and
    $\theta\in\mathbb{Z}_8$ is chosen by the verifier, following the paper's
    angle convention. The picture focuses on these quantum messages;
    the rest of the classical interaction is left implicit.
  }
  \note<2>{
    \footnotesize
    Replace the quantum preparation and transmission with remote state
    preparation: a classical interaction that leaves the prover holding a
    quantum state correlated with the verifier's classical information.
    The lower arrow represents this interactive subprotocol, rather than a
    single classical message. This gives the overall prepare-and-send route
    to a fully classical verifier; we will discuss the required guarantees
    and our particular construction later.
  }
\end{frame}
